\documentclass[11pt]{memoir}
\usepackage[usenames,dvipsnames]{color}
\usepackage{amsmath}%
\usepackage{amsfonts}%
\usepackage{amssymb}%
\usepackage{amsthm}%
\usepackage{graphicx}
\usepackage{ dsfont }
\usepackage[hidelinks]{hyperref} %hyperlink
\usepackage{subcaption} %NON SO SE SERVE
\usepackage{enumitem} %enumerazioni con le i
\usepackage{booktabs} %tabelle più belle
\usepackage{placeins} %forza le immagini in posizione
\usepackage{mathabx}
\usepackage{epsfig}
\usepackage{caption} %caption fuori da ambienti float
\usepackage{appendix} %appendice per i codici
\usepackage{dsfont}
\usepackage{color}
\usepackage{listings}
\usepackage{xcolor} %per colorare il testo

\def\R{\mathbb{R}}
\def\Z{\mathbb{Z}}
\def\H{\mathbb{H}}
\def\A{\mathbb{A}d\mathbb{S}}
\newcommand\norm[1]{\lVert#1\rVert}


\definecolor{mygreen}{RGB}{28,172,0} % color values Red, Green, Blue
\definecolor{mylilas}{RGB}{170,55,241}
\usepackage{algorithm} %Parte per gli algoritmi
\usepackage{algpseudocode}
\makeatletter
\makeatother


\usepackage[utf8]{inputenc}
\usepackage[english]{babel}
\usepackage{mathrsfs}
\usepackage{array}
\usepackage{rotating}
\usepackage{multirow}
\usepackage{amsmath,amsthm}
\usepackage{mathtools}
\usepackage{stmaryrd}
\usepackage{tikz-cd}
\usepackage{csquotes}
\usepackage{lscape}
\tikzcdset{
    startcode/.style={
        start anchor={[xshift=0.35cm]south west}
        },
    endcode/.style={
        end anchor={[xshift=0.35cm]north west}
    }
}
\usepackage{tikz}
\usepackage{upgreek}
\usepackage{varwidth}
\usepackage{hyperref}
\usepackage{biblatex}
\usepackage[colorinlistoftodos]{todonotes}
\addbibresource{refs.bib}
\usetikzlibrary{arrows}
\newcommand{\encadrer}[1]{{%
  
  \fbox{%
    \begin{varwidth}{\dimexpr\linewidth-2\fboxsep-2\fboxrule}
      #1%
    \end{varwidth}%
  }
  \par}%
}

\theoremstyle{plain}
\newtheorem{theorem}{Teorema}[chapter] % reset theorem numbering for each chapter
\theoremstyle{definition}
\newtheorem{definition}[theorem]{Definizione} % definition numbers are dependent on theorem numbers
\newtheorem{lemma}[theorem]{Lemma} % same for example numbers
\theoremstyle{plain}
\newtheorem{proposition}[theorem]{Proposizione}
\newtheorem{observation}[theorem]{Osservazione}
\renewcommand*{\proofname}{Dimostrazione}
\newcommand\restr[2]{\ensuremath{\left.#1\right|_{#2}}}
\newtheorem{corollary}[theorem]{Corollario}

\chapterstyle{veelo}
\checkandfixthelayout
\setcounter{tocdepth}{2}


\begin{document}
\advance\vsize by 2.5cm % Advance page height
\advance\voffset by -1.5cm % Shift top margin


\begin{titlingpage}


\advance\hsize by 2.7cm % Advance page height
\advance\hoffset by -1.35cm % Shift top margin


%%\aliaspagestyle{titlingpage}{empty}

%\newgeometry{margin=2.5cm}

\calccentering{\unitlength}
\begin{adjustwidth*}{\unitlength}{-\unitlength}


\thispagestyle{empty}
\begin{center}
\large
\textsc{\Large Università di Pisa\\}
\vspace{0.6cm}
\includegraphics[width=8cm]{cherubino.eps}

\vspace{0.8cm}
\textsc{\Large{Corso di laurea in Matematica}}\\



\vspace{2.5cm}


{\LARGE\textbf{Thurston Earthquake Theorem via Anti-de Sitter Geometry}}
\\[2.0cm]

\textsc{Tesi di Laurea Magistrale \\[0.2cm] in Matematica}\\
\vspace{0.5cm}

\begin{center}

\textsc{\Large Anno Accademico 2023/2024}
\end{center}

\vspace{0.2cm}

%\begin{center}
%\textsc{{\LARGE Candidato}}\\
%\textbf{{\small Tizio Caio}}
%\end{center}
%
%\vspace{0.4cm}
%
%\begin{center}
%\textsc{{\LARGE Relatore}}\\
%\textbf{{\small Sempronio Mevio}} \\
%{\small Università di Qualche Posto}
%\end{center}

\begin{center}
\makebox[\textwidth]{
\begin{tabular}{l p{5.3cm} r}
\textsc{{\LARGE Candidato}} & & \textsc{{\LARGE Relatore}} \\
\textbf{{\small SMC}} & & \textbf{{\small AT}} \\
  & & {\small Università di Pisa} \\
\end{tabular}}
\end{center}




\vspace{3.75cm}


\textsc{Anno Accademico 2023 - 2024}\\[0.2cm]


\end{center}

\end{adjustwidth*}


\advance\hsize by -2.7cm % Advance page height
\advance\hoffset by 1.35cm % Shift top margin


\end{titlingpage}


\advance\vsize by -2.5cm % Return old margings and page height
\advance\voffset by 1.5cm % Return old margings and page height



%\par\vfill\break % Break the page with different margins



%\restoregeometry

\par\vfill\break

\thispagestyle{empty}
\vspace*{\stretch{1}}
\begin{flushright}
%\itshape A Franca ed Emilia.

\end{flushright}
\vspace{\stretch{2}}
\cleardoublepage

\clearpage

\clearpage

\tableofcontents
\clearpage

\thispagestyle{empty}
\listoffigures

\thispagestyle{empty}
%\include{legenda}

%\include{ringraziamenti}
%\setcounter{chapter}{-1}
%\include{abstract}
\chapter{Lorentzian geometry}

In what follows we will introduce Lorentzian manifold of constant curvature and shows that, as in the Riemannian case, two manifolds of constant sectional curvature $K$ are locally isometric. We will then define what we mean by a manifold with maximal isometry group as they provide models of manifold with constant curvature: if we have $M$ a Lorentzian manifold with costant sectional curvature $K$ and maximal isometry group, then any Lorentzian manifold with constant sectional curvature $K$ carries a natural (Isom($M$), $M$)-atlas made of local isometries. \\
Simply connected space have maximal isometry group, but in general the converse is false. In particular we will focus on the case $K=-1$ and in that case it will be convenient to use non-simply connected models. 

\subsection{Basic Definitions}
By a \textit{Lorentzian metric} on a $n+1$-manifold we mean a non degenerate $2$-tensor $g$ of signature $(n,1)$. A \textit{Lorentzian manifold} is a connected manifold $M$ equipped with a Lorentzian metric $g$.\\
In a Lorentzian manifold $M$ we say that a non-zero vector $v \in TM$ is \textit{time-like} if $g(v,v)<0$, \textit{space-like} if $g(v,v)>0$ and \textit{light-like} if $g(v,v)=0$. More generally, we say that a linear subspace $V \subset T_{x}M$ is \textit{spacelike, lightlike, timelike} if the restricion of $g_x$ to $V$ is positive definite, degenerate or negative definite respectively.\\
The set of lightlike vectors, together with the null vector, disconnects $T_{x}M$ into 3 regions: two convex open cones formed by timelike vectors, one opposite to the other, and the region of spacelike vectors. 

\textcolor{red}{Inserire figura del cono di Minkowski?}

As a consequence the set of timelike in the total space $TM$ is either connected or is made by two connected components. In the latter case $M$ is said \textit{time-orientable}, and a time orientation is the choice of one of those components. Vectors is the chosen component are said     \textit{future-directed}, vectors in the other component are said \textit{past-directed}\\
A differentiable curve is said \textit{timelike, spacelike, lightlike} if its tangent vector at every point is timelike, spacelike, lightlike respectively. The curve is said to be \textit{causal} if the tangent vector is either timelike or lightlike. \\ Given a point $x$ in a time-oriented Lorentzian manifold $M$, the \textit{future} of $x$ is the set $I^+(x)$ of points which are connected to $x$ by a future-directed causal curve. The \textit{past} of $x$, $I^-(x)$, is defined in an analogus way for past-directed causal curves. 
As in the Riemannian setting, on a Lorentzian manifold $M$ there is a unique linear connection $\nabla$ which is symmetric and compatible with the Lorenzian metric $g$. We refer to it as the \textit{Levi-Civita connection} of $(M,g)$.\\ Following the analogy with Riemannian geometry the Levi-Civita connection determines the Riemann curvature tensor defined by: 
\[
    R(u,v)w=\nabla_u\nabla_v w-\nabla_v\nabla_u w-\nabla_{[u,v]}w
\]  

We then say that a Lorentzian manifold has constant sectional curvature $K$ if: 
\begin{equation}\label{sectionalcurvature}
    g(R(u,v)v,u)=K(g(u,u)g(v,v)-g(u,v)^2) 
\end{equation}
    


for every pair of vectors $u,v \in T_{x}M$ and every $x\in M$. Even tough the definition is analogus to the one given in the Riemannan realm we recall that in the Lorentzian setting the sectiona curvature can be defined only for planes in $T_{x}M$ where $g$ is non-degenerate. 

\subsection{Maximal isometry group and geodesic completeness} 

\begin{lemma}\label{isometrie} Let $M$ and $N$ be Lorentzian manifolds of constant curvature K. Then every linear isometry $L:T_{x}M\to T_yN$ extends to an isometry $f:U\to V$ where U, V are open neighbourhoods of $x,y$ respectively. Any two extensions $f:U\to V$ and $f^{\prime}:U^{\prime} \to V^{\prime} $ of $L$ coincide on $U\cap U^{\prime} $ Moreover $L$ extends to a local isometry $f:M\to N$ provided that $M$ is simply connected and $N$ is geodesically complete. 
\end{lemma}
\begin{proof}
    Cartan-Ambrose-Hicks thm.
\end{proof}
As a consequence of the aforementioned lemma we have: 
\begin{corollary}\label{122}
    Let $M$ and $N$ be simply connected, geodesically complete Lorentian manifolds of constant curvature $K$. Then any linear isometry $L:T_xM\to T_yN$ extends to a global isometry $f:M\to N.$   
\end{corollary}

In particular, there is a unique simply connected geodesically complete Lorentzian manifold of constant curvature $L$ up to isometries. For instance for $K=0$ a model is the Minkowski space $\R^{n,1}$, that is $\R^{n+1}$ provided with the standard metric: 
\[
    g=dx_1^{2}+\dots+dx_n^{2}-dx_{n+1}^2.   
\]

Another consequence of \ref{isometrie} is that, fixing a point $x_0 \in M$, the set of isometris of $M$, which we will denote by Isom($M$) can be realized as a subset of ISO($T_{x_0}M, TM$) namely the fiber bundle over $M$ whose fiber over $x\in M$ is the space of linear isometries of $T_{x_0}M$ into $T_{x}M$. \todo{Not clear what is going on here, he argues about some Lie structure.}

\begin{definition}
    A Lorentzian manifold $M$ has \textit{maximal isometry group} if the action of Isom($M$) is transitive and, for every point $x \in M$, every linear isometry $L:T_{x}M\to T_xM$ extends to an isometry of $M$. 
\end{definition}

Equivalently $M$ has maximal isometry group is the above inclusion of Isom($M$) into ISO($T_{x_0}M, TM$) is a bijection. Hence is $M$ has maximal isometry group, then the dimension of the isotropy group is maximal. 
From \ref{122}, every simply connected Lorentzian manifold $M$ has maximal isometry group if it has constant sectional curvature and is geodesically complete. The converse holds even without the simply connectedness assumption. Namely: 

\begin{lemma}\label{maximalisometry}
    If $M$ is a Lorentzian manifold with maximal isometry group, then $M$ has constant sectional curvature and is geodesically complete.
\end{lemma}

\begin{proof}
    Fix a point $x\in M$. As the identity component of $O(T_x M)\simeq O(n,1)$ acts transitebly on spacelike planes, tehre exists a constant $K$ such that \ref{sectionalcurvature} holds for every pair $(u,v)$ of vector tangent at $x$ which generate a spacelike plane. Now, for every point $x\in M$ both sides of equations \ref{sectionalcurvature} are polynomial in $u,v \in T_xM$. Since the set of pairs $(u,v)$ which generate spacelike planes is open in $T_{x}M\times T_{x}M$, the equation \ref{sectionalcurvature} must holds for every pair of vectors $u,v \in T_xM$.\\ Since Isom($M$) acts transitively on $M$, it follows that $M$ has constant sectional curvature $K$.\\
    Let us now show that the manifold is geodesically complete. Suppose $\gamma$ is a parametrized geodesic with $\gamma(0)=x$ and $\gamma^{\prime} (0)=v\in T_xM,$ which is definided for a finite maximal time $T>0.$ Let $T_0=T-\epsilon>0.$ By assumption one can find an isometry $f:M\to M$ such that $f(x)=\gamma(T_0)$ and $df_x(v)=\gamma^{\prime}(T_0).$ Then $t\to f\circ\gamma(t-T_0)$ is a parametrized geodesic which provides a continuation of $\gamma$ beyond $T$, a contradiction. 
\end{proof}

We conclude this preliminaries with a result of classification which will be useful in the following and also explains our interest in spaces with maximal isometry group.
\begin{proposition}\label{classification}
    Let $M_k$ be a simply connected Lorentzian manifold of constant sectional curvature $K$ with maximal isometry group, and let $M$ be a Lorentizian manifol of constant sectional curvature $K$. Then: 
    \begin{itemize}
        \item M is geodesically complete if and only if there is a local isometry $p:M_K\to M$ which is a universal covering.
        \item $M$ has maximal isometry group if and only if $\text(Aut)(p:M_k\to M)$ is normal in Isom($M_k$). 
    \end{itemize}
\end{proposition}

\begin{proof}
    Suppose $M$ is geodesically complete, then lifting the metric to the universal cover $\tilde{M}$ one gets a simply connected geodesically complete Lorenztian manifold of constant sectional curvature $K$ which by corollary \ref{122} is isometric to $M_K$. The covering map $p:_k\to M$ is then a local isometry by construction. The converse is straightforward. \\   
    Now let $\Gamma=\text{Aut}(p:M_K\to M),$ which is a discrete subgroup of Isom($M_K$). Thus $M$ is obtained as the quotient $M=M_{K}/\Gamma$,where $\Gamma$ acts freely and properly discontinuosly on $M_K$. The isometry group of $M$ is isomorphic to $N(\Gamma)/\Gamma$, where by $N(\Gamma)$ we denote the normalizer of $\Gamma$ in Isom($M_K$). The isomorphism is based on the observation that any isometry of $\tilde{M}$ which normalies $\Gamma$ descends to an isometry of $M$, and conversly the lifting of any isometry of $M$ must be in $N(\Gamma).$\\
    Hence the condition that $M$ has maximal isometry group is equivalent to the condition that every element $f$ of $\text{Isom}(M_K)$ descends in the quotient to an isometry of $M$. This is in turn equivalent to the condition that $f\Gamma f^{-1}=\Gamma$ for every $f\in \text{Isom}(M_K)$, which is the same as saying that $\Gamma$ is normal in $\text{Isom}(M_K)$.
\end{proof}

\begin{observation}
Since $\Gamma=\text{Aut}(p:M_K\to M)$ is discrete, being normal in $\text{Isom}(M_K)$ implies that elements of $\Gamma$ commute with the elements of the identity components of $\text{Isom}(M_K)$. This remark suggest that there are usually not many Lorenztian manifolds of constant sectional curvature with maximal isometry group.
\end{observation}

Finally, any isometry between connected open subsets of a Lorenztian manifold $M$ with maximal isometry group extends to a global isometry. In particular if $M_K$ is a Lorenziant manifold of constant sectional curvature $K$ with maximal isometry group, then any Lorentzian manifold $M$ of constant sectional curvature $K$ admits a natural $(\text{Isom}(M_K), M_K)$-structure whose charts are isometries between open subsets of $M$ and open subsets of $M_K$. We will refer to Lorentzian manifolds of constant sectional curvature $K$ with maximal isometry group as \textit{models} of constant sectional curvature $K$ (following Klein's terminology).  
\chapter{Model of Anti-de Sitter (n+1)-space}
The aim of this chapter is to construct models of Lorentzian manifolds with constant sectional curvature -1 and maximal isometry group in any dimension. We are also interested in \textit{stressing} the analogies of this manifolds with the models of hyperbolic space in the Riemannian setting. We will show that hyperbolic space is \textit{naturally} embedded in Anti-de Sitter space, and we will later develop this topic in our study of Earthquakes theory. After introducing the models, we are interested in studying the geometry of such manifolds: we will introduce a conformal boundary for Anti-de Sitter space and we will characterize geodesics and totally geodesic subspaces. We will also introduce the notion of polarity in Anti-de Sitter space (in same sense the \textit{correct} Lorentzian correspondence between points and hyperplane) and study its properties. 

\section{The quadric model}
We want to introduce the analogue of the hyperboloid model of hyperbolic space. Denote by $\R^{n,2}$ the real vector space $\R^{n+2}$ equipped with the quadratic form 
\[
    q_{n,2}(x)=x_1^{2}+\dots+x_n^{2}-x_{n+1}^{2}-x_{n+2}^2.   
\]
and by $<v,w>_{n,2}$ the associated symmetric form. Finally, let $O(n,2)$ be the group of linear transformation of $\R^{n+2}$ that preserve $q_{n,2}.$
Then we define: 
\[
    \H^{n,1}=\{x\in \R^{n,2}\;|\;q_{n,2}(x)=-1\}.
\]

It is immediate to check that $\H^{n,1},$ as the pre-image of a regular value of $q_{n,2}$, is a smooth connected submanifold of $\R^{n,2}$ of dimension $n+1$. The tangent space $T_x\H^{n+1}$, regarded as a subspace of $\R^{n+2}$ coincides with the orthogonal space $x^{\perp}=\{y\in \R^{n+2}\;|\;\langle x,y\rangle_{n,2}=0\}.$ A simple signature argument (and $q_{n,2}(x)=-1$ for every $x\in\H^{n,1}$) shows that the restriction of the symmetric form $\langle .,. \rangle_{n,2}$ to $T\H^{n+1}$ has Lorentzian signature, hence it makes $\H^{n,1}$ a Lorentzian manifold. We remark that this model is the analogue of the hyperboloid model of hyperbolic space, in fact $\H^n$ is isometrically embedded in $\H^{n,1}$ as the submanifold defined by $x_{n+2}=0$, $x_{n+1}>0$. \\
The natural action of $O(n,2)$ on $\R^{n,2}$ preserves $\H^{n,1}$, and in facts $O(n,2)$ acts by isometries on $\H^{n,1}$. We remark that $O(n,2)$ acts transitively on $\H^{n,1}$ and that the stabilizer of a point $x$ acts transitively on the space of orthonormal bases of $T_x\H^{n,1}$. Hence $\H^{n,1}$ has maximal isometry group and Isom($\H^{n,1})\simeq O(n,2)$.\\
By Lemma \ref{maximalisometry}, $\H^{n,1}$ has constant sectional curvature. Let us now check that the sectional curvature is negative (actually, $K=-1$). For this purpose, observe that the normal line in $\R^{n,2}$ to $\H^{n,1}$ at $x$ is identified with the line generated by $x$ itself. It follows that, if $v,w$ are tangent vector fields along $\H^{n,1}$ we have the orthogonal decomposition: 
\[
    (D_{v}w)(x)=(\nabla_{v}w)(x)+\langle v,w\rangle x, 
\]

where $D$ is the flat connection of $\R^{n+2}$ and $\nabla$ is the Levi-Civita connection of $\H^{n,1}$. Using the flatness of $D$ we get: 
\[
    R(u,v)w=\langle u,w\rangle v-\langle v,w\rangle u,
\]

so that 
\[
    \langle R(u,v)v,u \rangle =-(\langle u,u\rangle\langle v,v\rangle-\langle v,u\rangle^2), 
\]

and this shows that $\H^{n,1}$ has constant sectional curvature $-1$. We also remark that $\H^{n,1}$ is not simply connected, being homeomorphic to $\R^{n}\times \mathbb{S}^1$. 

\section{The ``Klein model" and it's boundary.}
Let us introduce a projective model, also knows as the ``Klein model", for Anti-de Sitter geometry. Let us define: 
\[
    \A^{n,1}=\H^{n,1}/\{\pm \text{Id}\}.
\]

Since $\{\pm \text{Id}\}$ is the center of $O(n,2)$ (hence normal), $\A^{n,1}$ (when endowed with the Lorentzian metric induced by the quotient) has maximal isometry group by Proposition \ref{classification} and is therefore a model of constant sectional curvature $-1$. It can also be shown that the center of the isometry group of the Klein model is trivial, hence $\A^{n,1}$ is the \textit{minimal} model of AdS geometry, in the sense that any other model is a covering of $\A^{n,1}$.\\
By definition $\A^{n,1}$ is naturally identified with a subspace of real projective space $\R \text{P}^{n+1}$, more explicitly with the subset of timelike directions of $\R^{n,2}$: 
\[
    \A^{n,1}=\{[x]\in \R \text{P}^{n+1}\;|\;q_{n,2}(x)<0\}.
\]
Like in hyperbolic geometry, the boundary of $\A^{n,1}$ in projective space is a quadric of signature $(n,2)$, that is the projectivization of the of lightlike vectors in $\R^{n,2}$. We denote this quadric by $\partial \A^{n,1}=\{[x]\in \R \text{P}^{n+1}\;|\;q_{n,2}(x)=0\}$.\\ We observe that isometries of $\A^{n,1}$ induce projective transformations which preserve $\partial \A^{n,1}$.  \\

\textit{The conformal Lorentzian structure of the boundary.}
We want to continue to develop analogy with hyperbolic geometry and equip $\partial \A^{n,1}$ with a conformal Lorentzian structure that extends the conformal Lorentzian structure of $\A^{n,1}$. \\
A point $\ell \in\R\text{P}^{n+1}$ is identified with $\text{Span}(x)$ for some $x\in \R^{n,2},$ and the tangent space of real projective space has the canonical identification 
\[
    T_{\ell}\R\text{P}^{n+1}\simeq \text{Hom}(\ell,\R^{n+2}/\ell). 
\]

Now if $\ell$ is timelike we can identify $\R^{n+2}/\ell$ with $\ell^{\perp}.$ For any given local section $\sigma:\A^{n,1}\to \R^{n,2}$ of the projection $\R^{n,2}\to \A^{n,1}$, we can define a Lorentzian metric on $T\A^{n,1}$ setting:

\[
    \langle \langle f,g \rangle \rangle_{\sigma}=\langle f((\sigma[x])),g(\sigma[x])\rangle_{n,2}. 
\]

for $f,g\in T_{x}\A^{n,1}\simeq \text{Hom}(\ell,\ell^\perp$). If the section $\sigma$ has image in $\H^{n,1}$, then the aforementioned metric coincides with the pull-back of the metric over $\R^{n,2}$, since the differential of $\sigma $ identifies $T_{[x]}\A^{n,1}=T_x\H^{n,1}=x^\perp$. For a general section the identification does not hold, but we can recover a conformal metric via the formula: 
\[
    \langle \langle f,g \rangle \rangle_{\lambda\sigma}=\lambda^2\langle \langle f,g \rangle \rangle_{\sigma}.
\] for any function $\lambda$.\\
We consider now the case that $\ell=\text{Span}(x)$ is lightlike, i.e. $q_{n,2}(x)=0$. In this case we can not induce any natural metric on $\R^{n,2}/\ell$. However, we let 
\[
    \mathbb{L}=\{x\in \R^{n,2}\;|\;q_{n,2}(x)=0\}
\]
be the space of lightlike vectors, then $T_x\mathbb{L}$ is precisely $\ell^{\perp}$ and contains $\ell$ itself. We have a canonical identification: $T_\ell\partial\A^{n,1}\simeq\text{Hom}(\ell,\ell^{\perp }/\ell).$ The bilinear for of $\R^{n,2},$ when restricted to $\ell^{\perp},$ induces a non degenerate bilinear form (of signature $(n-1,1)$) on $\ell^\perp/\ell.$ We will denote such a restriction as $\langle v,w\rangle_{\ell^{\perp}/\ell}.$ \\
We can now define a metric on $\partial\A^{n,1}$ for any section $\sigma:\partial\A^{n,1}\to \mathbb{L}$ of the canonical projection by the formula:
\begin{equation}\label{24}
    ((f,g))_{\sigma}=\langle f(\sigma[x]),g(\sigma[x])\rangle_{\ell^{\perp}/\ell},
\end{equation}
for all $f,g\in \text{Hom}(\ell, \ell^\perp/\ell).$ This metric can be viewed as the pull-back of the metric: 
\begin{equation}\label{25}
    ((f,g))_\sigma=\langle \sigma_*(f), \sigma_{\ast} (g)\rangle_{n,2}, 
\end{equation}
since the degenerate metric on $T_x\mathbb{L}=\ell^\perp$ is, by construction, the pull-back of the metric of $\ell^{\perp}/\ell$ by the projection along the degenerate direction $\ell$.\\ 
The relation valid for the metric on $\A^{n,1}$ also holds for the metric on $\partial\A^{n,1}$, that is: 
\begin{equation}\label{26}
    ((f,g))_{\lambda\sigma}=\lambda^2((f,g))_{\sigma}
\end{equation} 
and therefore the induced conformal class over $T\partial\A^{n,1}$ is well defined and independent of the choice of $\sigma$ and equips the tangent space of the boundary with a conformal Lorentzian metric. To conclude let us show that this conformal metric is naturally the compactification of $\A^{n,1}$. Let $\sigma$ be a section, of the projection $\pi:\R^{n,2}\to \R\text{P}^{n+1},$ defined in a neighborhood $U$ of a point in $\partial\A^{n,1}$, by construction the metric $((\cdot,\cdot))_\sigma$ over $\partial\A^{n,1}\cap U$ is the limit of the conformal metric associated to $\sigma$ defined over $\A^{n,1}\cap U$. % This means that if $(p_n,v_n)$ is a sequence in $T\A^{n,1}$ that converges to $(p_{\infty}, v_\infty)$ then 
\begin{observation}\label{222}
Let's make some observation about the light cone in the case of $\partial\A^{n,1}$. If $[y]\in\partial\A^{n,1}$ Equation \refeq{25} implies that the lightlike vectors in $T_{[y]}\partial\A^{n,1}$ are exactly the projection of vectors $x\in \R^{n,2}$ such that $\langle x,y\rangle_{n,2}=0$ and $q_{n,2}(x)=0$. These vectors are such that $\text{Span}(x,y)$ are totally degenerate planes in $\R^{n,2}$, equivalently are projective lines contained in $\partial\A^{n,1}$. The light cone in $\partial\A^{n,1}$ through $[y]$ is therefore the union of all the projective lines through $[y]$ that are contained in $\partial\A^{n,1}$.

\end{observation}


\subsection{The Poincaré model for the universal cover.}
We have already observed that $\H^{n,1}$, and its double quotient $\A^{n,1}$, are not simply connected. We want to construct a simply connected model for AdS geometry. For this purpose we introduce the universal cover of $\H^{n,1}$ and $\A^{n,1}$.\\
Let $\H^n$ be the hyperboloid model of hyperbolic space. Then: 
\begin{equation}\label{ogcover}
    \pi(y,t)=(y_1,\dots,y_n,y_{n+1}\cos t,y_{n+2}\sin t)
\end{equation}


defines a map $\pi:\H^n\times \R\to \H^{n,1}$ which is a covering with deck transformations of the form $(y,t)\mapsto (y,t+2k\pi)$ for $k\in \Z$. We denote the covering space by $\AS^{n,1}$ and we observe that it is also the universal cover of $\A^{n,1}$, where the covering map it is just the composition of $\pi$ and the double quotient.\\
Pulling back the Lorentzian metric over $\AS^{n,1}$ we get a simply connected Lorentzian manifold of constant curvature -1. The metric on $\AS^{n,1}$ is a warped product of the form: 
\begin{equation}\label{metric}
     \pi^*g_{\H^{n,1}}=g_{\H^n}-y_{n+1}^{2}dt^2.
\end{equation}
   

Moreover $\AS^{n,1}$ has maximal isometry group, hence we have obtained a simply connected model for AdS geometry. More precisely we have a central extension, that is a (non split) short exact sequence:
\[
    0\to \Z\to \textnormal{Isom}(\AS^{n,1})\to O(n,2)\to 1.
\]
It is sometimes convenient to express the metric \ref{metric} using the Poincaré model of hyperbolic space. Recall that the disk model of the hyperbolic space is the unit disk $\D^n$ endowed with the conformal metric: $\frac{4}{(1-r^2)^2}\sum dx_i^2$, where $r^2=\vert x\vert^2$. In our setting the isometry between the disk and the hyperboloid model is given by: 
\begin{equation}
    (x_1,\dots, x_n)\mapsto (y_1=\frac{2x_1}{1-r^2},\dots, y_n=\frac{2x_n}{1-r^2},y_{n+1}=\frac{1+r^2}{1-r^2}) 
\end{equation}
The Poincaré model of AdS geometry is then the cylinder $\D^n\times \R$ endowed with the metric 
\begin{equation}\label{metricdisk}
    \frac{4}{(1-r^2)^2}(dx_1^2+\dots+dx_n^2)-(\frac{1+r^2}{1-r^2})^{2}dt^2 
\end{equation}

It follows from the definition that each slice $\{t=c\}$ is a totally geodesic copy of $\H^n$. The metric defined in \refeq{metricdisk} also shows that the vector field $\partial/\partial t$ is a timelike non vanishing vector field on $\AS^{n,1}$, giving to $\AS^{n,1}$ the structure of a time-orientable manifold. Any choice of time orientation is preserved by the action of the deck transformation of the covering from the Poincaré to the Klein model, hence both $\H^{n,1}$ and $\A$ are time-orientable. \\
\begin{observation}
    Exception made for the central line passing through $x_1=\dots=x_n=t=0,$ vertical lines are not geodesic for the metric defined in \refeq{metricdisk}.  
\end{observation}

\textit{The conformal metric of the boundary.} As we did with $\A^{n,1}$ we would like to give a conformal Lorentzian structure to the boundary of $\AS^{n,1}$. The construction actually works for any covering space of $\A^{n,1}.$ The covering map (we are now considering the projective model of $\H^n$)
\[
    \pi^{\prime} ([y_1:\dots:y_{n}:y_{n+1}],t)=[y_1:\dots:y_n:y_{n+1}\cos t:y_{n+1}\sin t]
\]

extends to a map $\pi^{\prime}:\H^n\cup \partial\H^n\times \R\to \A^{n,1}\cup \partial\A^{n,1}$. In order to compute the conformal structure of the boundary, we consider the map $\tau: \H^n\times \R\to\R^{n+2}$ defined by:
\[
    \tau([y_1:\dots:y_n:y_{n+1}],t)=(y_1/y_{n+1},\dots,y_n/y_{n+1},\cos t,\sin t)
\]  
 that extends to the boundary and induces a (local) section of the projection $\R^{n,2}\to \A^{n,1}$. In fact, let $\eta$ be the generator of the group of deck transformation of the covering $\pi^{\prime}:\AS^{n,1}   \to \A^{n,1}$, then $\tau$ has the equivariance $\tau\circ\eta^i=(-1)^i\tau.$  % Using Equation \refeq{26}, the conformal Lorentzian metric on $\partial\AS^{n,1}$ induced by $\sigma$ by means of Equation \refeq{24} is the pull-back of a Lorentzian metric compatible with the natural conformal structure of the boundary $\partial\A^{2,1}$.  
\subsection{Geodesic.}

We have presented our various models as manifolds and now we would like to improve our knowledge of their geometry. As always we start by studying and characterizing geodesics.\\ 

\textit{In the quadric model} Let us start with the exponential map in the hyperboloid model. Given a point $x\in \H^{n,1}$ and a vector $v\in T_x\H^{n,1}$, we want to determine the geodesic through $x$ with speed $v$. We will distinguish several cases according to the sign of $ q_{n,2}(v)$. If $v$ is lightlike, then: 
\[
    \gamma(t)=x+tv
\] is a geodesic of $\R^{n,2}$ and is contained in $\H^{n,1}$, hence $\gamma$ is a geodesic of $\H^{n,1}$. If $v$ is etiher timelike or spacelike, we claim that the geodesic $\gamma(t)=\exp_x(tv)$ is contained in the linear plane $W=\text{Span}(x,v).$ In fact, the linear transformation $T$ that fixes pointwise $W$ and whose restriction to $W^\perp$ is $-\text{Id}_{W^\perp}$ is in $O(n,2)$. By the uniqueness of the geodesic, $T\circ\gamma=\gamma$ hence $\gamma$ is contained in $\H^{n,1}\cap W$. We can easily derive the expressions\\
\begin{equation}
    \gamma(t)=\cosh(t)x+\sinh(t)v
\end{equation}

if $q_{n,2}(v)=1$ and 
\begin{equation}\label{212}
    \gamma(t)=\cos(t)x+\sin(t)v
\end{equation}
if $q_{n,2}(v)=-1$.\\

\textit{In the Klein model.} In analogy with the Riemannian case, in the Klein model $\A^{n,1}$ geodesics are intersections of projective lines with the domain of $\A^{n,1}\subset\R \text{P}^{n+1}.$ From what we have already said: 
\begin{itemize}
    \item Timelike geodesics correspond to projective lines that are entirely contained in $\A^{n,1}$, are closed non-trivial loops and have length $\pi.$ 
    \item Spacelike geodesics correspond to lines that meet $\partial \A^{n,1}$ transversally in two points. They have infinite length.
    \item Lightlike geodesics correspond to lines tangent to $\partial\A^{n,1}$.   
\end{itemize}

In particular the light cone through a point $[x]\in \A$ coincides with the cone of lines through $[x]$ tangent to $\partial \A$.\\ For instance in the affine chart $\mathbb{A}_{n+2}\neq\{x_{n+2}=0\},$ where in coordinates $(y_1, \dots, y_{n+1})=(x_1/x_{n+2},\dots, x_{n+1}/x_{n+2})$, the intersection $\A\cap\mathbb{A}_{n+2}$ is the interior of a one sheeted hyperboloid, that is: 
\[
    \A\cap\mathbb{A}_{n+2}=\{y_1^2+\dots+y_n^2-y_{n+1}^2<1\},
\]
 while the boundary it's the one-sheeted hyperboloid itself: 
 \[
    \partial\A\cap\mathbb{A}_{n+2}=\{y_1^2+\dots+y_n^2-y_{n+1}^2<1\}.
\]
In an affine chart, timelike geodesic corresponds to affine lines which are entirely contained in the Anti de Sitter space, and which are not asymptotic to its boundary; lightlike geodesics are tangent to the one sheeted hyperboloid, or are asymptotic to it (tangent at infinity).

\begin{observation} An important observation concerns the space of timelike geodesics. Any timelike line is the projectivisation of a negative definite plane. As Isom$(\A^{n,1})\simeq \text{PO}(n,2)$ acts transitively on the space of timelike lines, and since the stabilizer of a timelike line is the group $\text{P}(O(n)\times O(2))$ which is the maximal compact subgroup of $\text{PO}(n,2)$, the space of timelike geodesics of $\A^{n,1}$ is naturally identified with the Riemannian symmetric space of PO$(n,2)$.
\end{observation}

\noindent\textit{Totally geodesic subspaces} Totally geodesic subspaces of $\A^{n,1}$ of dimension $k$ are obtained as the intersection of $\A^{n,1}$ with the projectivisation $P(W)$ of a linear subspace $W$ of $\R^{n,2}$ of dimension $k+1.$ The negative index of $W$ can be $1, 2,$ for otherwise the intersection $\A^{n,1}\cap P(W)$ would be empty. We have several cases: 
\begin{itemize}
    \item If $W$ has signature (k-1,2), then $P(W)\cap \A^{n,1}$ is isometric to $\A^{k-1,1}$. 
    \item If $W$ has signature $(k-2, 1)$ , then it is a copy of Minkowski space $\R^{k-2,1}$, hence $P(W)\cap \A^{n,1}$ is a copy of the Klein model of hyperbolic space. 
    \item If $W$ is degenerate, then $P(W)\cap \A^{n,1},$ is a lightlike subspace foliated by lightlike geodesics tangent to the same point of $\partial\A^{n,1}$. 
\end{itemize}  

A particular case of the last point is when $W$ is degenerate and $\text{dim}(W)=n+1.$ Then $P(W)\cap \A^{n,1}$ is a projective hyperplane tangent to $\partial\A$ at a point $[x]$ and $P(W)\cap \partial\A^{n,1}$ is the lightlike cone of $\partial\A^{n,1}$ through $[x]$. \todo{Check remark 2.2.2} 

\textit{In the universal cover.} In the universal cover $\AS^{n,1},$ geodesics are just the lifts of geodesics in $\A^{n,1}$ or $\H^{n,1}$. Hence every spacelike or lightlike geodesic in $\A^{n,1}$ and $\H^{n,1}$, which is topologically a line, has a countable number of lifts to $\AS^{n,1}$. Timelike geodesics in $\A^{n,1}$ and $\H^{n,1}$ are topologically circles and are in bijections with timelike geodesics in $\AS^{n,1},$ as the covering map restricted to a timelike geodesic, induces a covering map onto the circle. Using the Poincaré model we can give an explicit description of lightlike geodesic. In fact, in Lorentzian geometry not only the nature of a vector is conformally invariant but also unparametrized lightlike are a conformal invariant: 
\begin{theorem}\label{ConformalMetric} If two Lorentzian metrics $g,g^{\prime} $ on a manifold $M$ are conformal, then they have the same unparametrized lightlike geodesics.
\end{theorem}
\begin{proof}
    [GHL04, Proposition 2.131]
\end{proof}
Because of Theorem \ref{ConformalMetric} we can replace the Poincaré metric by the conformally equivalent metric given by:
\begin{equation}\label{emispherical}
    \frac{4}{(1+r^2)^2}(dx_1^2+\dots+dx_n^2)-dt^2
\end{equation} 
Now we observe that the first term in Equation \ref{emispherical} is exactly the form of the spherical metric on a hemisphere, pulled-back to the unit disk by the stereographic projection. We will call such a metric hemispherical and will denote it by $g_{\mathbb{S}^n}$. Notice that the boundary of $\partial\D$ is an equator for the hemispherical metric, and in fact it's the only equator completely contained in $(\D\cup\partial\D,g_{\mathbb{S}^n})$, a justification to the fact that we will refer to it as \textit{the} equator.\\
As a consequence, unparametrized lightlike geodesic of $\AS^{n,1}$, going through a point $(p_0,t_0)$ are characterized by the conditions that they are mapped to spherical geodesic under the vertical projection $(p,t)\to p$ and moreover: 
\[
    t-t_0=d_{\mathbb{S}^n}(p,p_0)
\] on the geodesic. In particular, these lightlike geodesic meet the boundary of $\AS^{n,1}$ at the point that satisfies the above conditions such that $p$ is on the equator of the hemisphere: as an example, if $p_0$ is the center of the hemisphere, then the points at infinity of the lightcone over $(p_0,t_0)$ are the horizontal slice $t=t_{0}+\pi/2.$ This sphere is also the boundary of a hyperplane dual to $(p_0,t_0)$, in a sense that we will explain in the following section. 

\section{Polarity in Ads}\label{polarsec}
The quadratic form $q_{n,2}$ induces a polarity on the projective space $\R\text{P}^{n+1}$, explicitly the correspondence associates to a projective subspace $P(W)$ the subspace $P(W^\perp)$. In particular, we have an induced duality between spacelike totally geodesic subspaces of $\A^{n,1}$ where the dual of a spacelike $k-$dimensional subspaces is an $n-k+1$ subspace. \\
For instance, if we consider the dual of a point $[x]\in\A^{n,1}$ it will be an $n-$dimensional spacelike hyperplane $P_{[x]}=\text{P}^\perp.$ Projectively $P_{[x]}$ is characterized as the hyperplane spanned by the intersection of $\partial\A^{n,1}$ \todo{n o n-1} with the lightcone from $[x]$. More geometrically, it can be checked that $P_{[x]}$ is the set of antipodal point to $[x]$ along timelike geodesics through $[x].$ Also, every timelike geodesic through $[x]$ meets $P_{[x]}$ orthogonally at time $\pi/2.$ Conversely, given a totally geodesic spacelike hyperplane $H$, all the timelike geodesic that meet $H$ orthogonally intersect in a single point, which is the dual point of $H$.

\subsection*{In the quadric model} We would like to \textit{lift} the idea of duality to coverings of $\A^{n,1}.$ Observe that in $\H^{n,1}$ there are two dual planes associated to any point $x$, namely the sets: 
\[
    P_x^\pm=\{\exp_x(\pm(\pi/2)v)\;|\;q_{n,2}(v)=-1,\; v\;\text{future-directed}\}.
\] 

Now the points $P_x^+$ and $P_x^-$ are antipodal and $P_{-x}^\pm=P_x^\mp$. The planes $P_x^\pm$ disconnect $\H^{n,1}$ in two regions $U_x$ and $U_{-x},$ where $U_x$ is the connected component of $x$. They can be characterized as: 
\[
    U_x=\{y\in\H^{n,1}\;|\;\langle x,y\rangle_{n,1}<0\}.
\]

Spacelike and lightlike geodesic through $x$ do not exit $U_x$, while all the timelike geodesics through $x$ meet orthogonally $P_x^\pm$ and all pass through the point $-x$. More precisely a point $y\neq x$ is connected to $x$: 
\begin{itemize}
    \item by a spacelike geodesic if and only if $\langle x,y\rangle_{n,1}<-1,$
    \item by a lightlike geodesic if and only if $\langle x,y\rangle_{n,1}=-1,$
    \item by a timelike geodesic if and only if $|\langle x,y\rangle_{n,1}|<1.$ 
\end{itemize}

An immediate consequence is that if $y$ is connected to $x$ by a spacelike geodesic, there is no geodesic joining $y$ to $-x$. Hence the exponential map of $\H^{n,1}$ is not surjective. But as any point $y\in\H^{n,1}$ can be connected through a geodesic either to $x$ or $-x$ hence the exponential over $\A^{2,1}$ is surjective.
\subsection{In the universal cover.}
Recall that $\AS^{n,1}\simeq \H^n\times\R$ with the group of deck transformations being isomorphic to $\Z$, where a generator acts by translation of $2\pi$ in the $\R$ factor. It follows that the preimage of a spacelike plane $P\subset\A^{n,1}$ is the disjoint union of spacelike planes $(P^k)_{k\in\Z},$ enumerated so that the generator $\eta$ of $\Z$ acts by sending $P^k$ to $P^{k+1}.$ \\
Each connected component of $\AS^{2,1}\setminus\bigcup_{k\in\Z} P^k$ is a fundamental domain for the action of deck transformation of the covering $\AS^{n,1}\to\A^{n,1}.$ Given a point $x$, let us apply the previous construction to the plane $P_x=P_{\pi^{\prime}(x)}$ which is the dual of the image $\pi^{\prime}(x)$ in $\A^{n,1}$, and let $V_x$ be the connected component which contains $x$. We will refer to $V_x$ as the \textit{Dirichlet domain} in $\AS^{2,1}$ centered at $x$, since the construction of $V_x$ is the analogue of a Dirichlet domain to our context. Then the restricted covering map: $\restr{\pi^{\prime}}{V_x}:V_x\to\A^{2,1}\setminus P_x$ is an isometry. Therefore lightlike and spacelike geodesic through $x$ are entirely contained in $V_x$.\\ 
\chapter{Anti-de Sitter space in dimension (2+1)}
We will now restrict our attention to dimension three Anti-de Sitter geometry, as it will be the specific model geometry of the manifolds of our next interest. In this chapter we will highlight the structure of Lie group that $\A^{2,1}$ has in this specific dimension and we will study its geometry using tools of Lie group theory. Most of the results are just a particular case of the theory developed in the previous chapter, but the Lie structure permits to give a more explicit description of the geometry of the phenomena we are interested in.\\

\section{The {PSL}$(2,\R)$ model.} 
The fundamental observation is the existence of a special model in dimension three which endows Anti-de Sitter space with a Lie group structure. To construct such a model we observe that $q=-$det is a quadratic form with signature (2,2) over the real vector space $\mathcal{M}(2,\R)$ (the signature is evident when we consider the basis consisting of elementary matrices). The associated bilinear form is expressed by the formula:
\begin{equation}\label{quadratic}
    \langle A,B\rangle=-\frac{1}{2}\text{tr}(A\cdot\text{adj}(B))
\end{equation}
for $A,B\in\mathcal{M}(2,\R)$, where adj denotes the adjugate matrix, namely: 
\[
    \text{adj}\Big(\begin{bmatrix}
        a & b \\
        c & d
    \end{bmatrix}\Big) = \begin{bmatrix}
        d & -b \\
        -c & a
    \end{bmatrix}
\]

hence, via Sylvester's theorem, there is an identification between $(\mathcal{M}(2,\R),q)$ and $(\R^{2,2}, q_{2,2}),$ unique up to composition by elements in $O(2,2)$. Under this isomorphism $\H^{2,1}$ is identified with the Lie group $\text{SL}(2,\R).$\\
Let us observe that $\text{SL}(2,\R)\times \text{SL}(2,\R)$ acts linearly on $\mathcal{M}(2,\R)$ by left and right multiplication:
$$(A,B)\cdot X=AXB^{-1}.$$

As a direct consequence of the Binet formula, the action preserves the quadratic form $q=$-det and thus induces a representation: 

\[ \rho:\text{SL}(2,\R)\times \text{SL}(2,\R)\to O(\mathcal{M}(2,\R),q). \]

Since the center of $\text{SL}(2,\R)$ is $\{\pm \text{Id}\},$ the kernel of $\rho$ is $K=\{(\text{Id},\text{Id}),(-\text{Id},-\text{Id})\},$ and by a dimensional argument (and connectedness of $\text{SL}(2,\R)$) it turns out that the image of the representation is the connected component of the identity: 
\[
    \text{Isom}_0(\H^{2,1})=\text{SO}_0(\mathcal{M}(2,\R),q)\simeq (\text{SL}(2,\R)\times \text{SL}(2,\R))/K
\]
    
Using this model, one has a natural identification between $\A^{2,1}$ and the Lie group $\PSL$, in such a way that: 
\begin{equation}\label{isomiso}
    \text{Isom}_0(\A^{2,1})\simeq \PSL\times\PSL
\end{equation}
    
acting by left and right multiplication on $\PSL$.\\ We are mostly interested in orientation-preserving and time-preserving notions that do not depend on a chosen orientation, nevertheless we will fix here an orientation and a time-orientation of $\A^{2,1}\simeq\PSL.$ As we are dealing with a Lie group it is sufficient to define on orientation of the Lie algebra, namely the tangent at the identity Id. We declare as (positive) oriented basis of $\mathfrak{sl}(2,\R):$ 
\begin{equation}\label{tangent}
V=\begin{pmatrix}
  0 & 1 \\ 1 & 0
\end{pmatrix}\;\;
W=\begin{pmatrix}
  1 & 0 \\ 0 & -1
\end{pmatrix}
\;\;
U=\begin{pmatrix}
  0 & -1 \\ 1 & 0
\end{pmatrix}
\end{equation}

The first two vectors $V,W$ are spacelike, while $U$ is timelike. $U$ is the tangent vector to the one-parameter group of elliptic isometries of $\H^2$ fixing $i\in\H^2$, parametrized by the angle of clockwise rotation; $V,W$ are vectors tangent to the one-parameter groups of loxodromic isometries fixing the geodesic with endpoints $(-1,1)$ and $(0,\infty)$ respectively. Time-orientation can also be inherited by the Lie algebra, we declare that $U$ is future-pointing timelike vector. \\

The stabilizer of the identity in Isom$_0(\A^{2,1})$ is the diagonal subgroup $\Delta<\PSL\times\PSL$. Under the obvious identification of $\PSL$ and $\Delta,$ the action of the identity stabilizer on the Lie algebra $\mathfrak{sl}(2,\R)=T_{\text{Id}}\PSL$ is the adjoint action of $\PSL$. A direct consequence of this construction is the bi-invariance of the quadratic form $q$. Indeed, denoting by $q_{\text{Id}}$ the restriction of $q$ to $T_{\text{Id}}\text{SL}(2,\R)$, a direct computation shows that $q_\text{Id}$ equals $(1/8)\kappa$ where $\kappa(X,Y)=4tr(XY)$ is the Killing form of $\mathfrak{sl}(2, \R)$.

\begin{observation}\label{311}
    The Lie algebra $\mathfrak{sl}(2,\R)$ equipped with the quadratic form $q_{\text{Id}}$ is then a copy of the 3-dimensional Minkowski space, hence the adjoint action yields a representation 
    \[
        \PSL\to O(\mathfrak{sl}(2,\R),q_{\text{Id}})
    \]

which in turn induces the well-known isomorphism: 
\[
    \text{SO}_0(2,1)\simeq\text{SO}_0(\mathfrak{sl}(2,\R), q_\text{Id})\simeq \PSL,
\] which is nothing but the restriction of the isomorphism of Equation \ref{isomiso} to the stabilizer of the identity in the left-hand side $\text{Isom}(\A^{2,1}),$ and to the diagonal subgroup $\Delta$ in the right-hand side $\PSL\times\PSL.$
\end{observation}

\section{The boundary of $\PSL$} 
From the aforementioned identification we obtain a new one between $\partial\A^{2,1}$ with the boundary of $\PSL$, inside $\text{P}(\mathcal{M}(2,\R))$ as the projectivization of the cone of rank 1 matrices. Therefore from now on we shall consider 
\[
    \partial\A^{2,1}=\{[X]\in\text{P}(\mathcal{M}(2,\R))|\;\text{rank}(X)=1\}.
\]
We endow $\overline{\A^{2,1}}\coloneqq \A^{2,1}\cup\partial\A^{2,1}$ with the topology induced by seeing both as subsets of the real projective space $\text{P}(\mathcal{M}(2,\R)).$ We want to observe that we have the following homeomorphism: 
\begin{align*}\label{delta}
    \delta:\partial\A^{2,1}&\rightarrow\R\text{P}^1\times\R\text{P}^1\\
    [X]&\mapsto (\text{Im}(X),\text{Ker}(X))
\end{align*}


where we are considering $\R\text{P}^1$ as the space of one-dimensional subspaces of $\R^2$. Since we have that $\text{Im}(AXB^{-1})=A\cdot\text{Im}(X)$ and $\text{Ker} (AXB^{-1})=B\cdot\text{Ker}(X),$ the map $\delta$ is equivariant with respect to the action of $\PSL\times\PSL,$ acting on $\partial\A^{2,1}$ as the natural extension of the group of isometries of $\A^{2,1}$ and on $\R\text{P}^1\times\R\text{P}^1$ by the obvious product action.\\
In this setting our choice of a time-orientation \textit{can be modified} according to the following Lemma:
\begin{lemma}\label{invertime}
    The inversion map $\iota[X]=[X]^{-1}$ is a time-reversing isometry of $\A^{2,1}$ which induces the homeomorphism $(x,y)\to(y,x)$ on the boundary $\partial\A^{2,1}.$
\end{lemma}
\begin{proof}
    \todo{intendono dire che $\iota[(A,B)(X)]=(\varphi(A,B))[\iota(X)]$} It follows from definition that $\iota$ is equivariant with respect to the isomorphism of $\PSL\times\PSL$ which switches left and right factors. As $\text{d}_\text{Id}\iota=-\text{Id}$ is a linear isometry, $\iota$ is an isometry, the differential being minus the identity also shows time-reversal. \\ 
    The second part of the statement can be checked via the following. For a $2\times2$ matrix the Cayley-Hamilton theorem implies the equality $(\det X)X^{-1}=(\text{tr}X)(\text{Id}-X),$ so that projectively we have $[X^{-1}]=[\text{tr}X\text{Id}-X].$ Hence $\iota$ extends to the transformation $[X]\to[\text{tr}X\text{Id}-X]$ on $\partial\A^{2,1}.$ Now if $X$ is a rank $1$ matrix, it is traceless if and only if $X^2=0$, hence $\text{Ker}(X)=\text{Im}(X).$ If $\text{tr}(X)\neq 0$, then $X$ is diagonalizable with eigenvalues $0$ and $\text{tr}(X)$. Moreover $\text{Ker}(X)$ and $\text{Im}(X)$ are the corresponding eigenspaces. It follows then that $\text{Ker}(\text{tr}X\text{Id}-X)=\text{Im}(X)$ and $\text{Im}(\text{tr}X\text{Id}-X)=\text{Ker}X$ 
\end{proof}


Consider the hyperbolic model of the upper half-plane $\H^2.$ $\R\text{P}^1$ correspond to the boundary at infinity $\partial\H^2$ and $\PSL$ is identified to $\text{Isom}_0(\H^2$). From this perspective we can consider $\partial\A^{2,1}$ as $\partial\H^2\times\partial\H^2$. We can then interpret the convergence to $\partial\A^{2,1}$ in this setting:

\begin{lemma}\label{convergenza}
    A sequence $[X_n]\in\A^{2,1}$ converges to $(x,y)\in\partial\A^{2,1}\simeq\T$ if and only if for every $p\in\H^2,$ $X_n(p)\to x$ and $X_n^{-1}(p)\to y.$ 
\end{lemma}
\begin{proof}
    $\PSL$ acts on $\H^2$ transitively via isometries, hence if the conditions holds for some $p$ then it holds for all $p\in\H^2$. Without loss of generality we can assume $p=i$ in the upper half-plane. Assuming $X_n$ converges projectively \todo{controllare se funziona} to a rank 1 matrix $X$, one checks immediately that $X(p)$ is in the projective class of $x=\text{Im}(X)$ as follows: let $X_n$ be the divergent sequence, then for any convergent sequence $v_n \in \mathbb{C}$ that does not accumulate on $\ker X$ we have that the limit of $X_n v_n$ is a non zero vector in $\text{Im}X$, so that 
    \[
        [X_n][v_n]=[X_n v_n]\rightarrow\text{Im}X\in\S.
    \]    Using this fact we have that $[A_n]\rightarrow (x,y)$ if and only if $[A_n]z\rightarrow x$ for $z\in\overline{\H^2}\setminus\{y\}$ and the convergence is uniform on compact subsets of $\overline{\H}^2\setminus\{y\}.$ \\
    The convergence of $X_n^{-1}(p)\to y$ is just a straightforward application of Lemma \ref{invertime} and an analogous argument to the one given for $X_n(p)$.
\end{proof}

In dimension three, $\partial\A^{2,1}$ is a doubly ruled quadric, which in an affine chart looks like \textcolor{red}{inserire cono}. We shall describe such rulings in a more geometric way. Given any $(x_0,y_0)\in\partial\A^{2,1}$, 
\begin{equation}\label{ruling}
        \lambda_{y_0}\coloneqq\{(x,y_0)|x\in\R\text{P}^1\}
\end{equation}

describes a projective line in $\R\text{P}^3$ which is contained in $\partial\A^{2,1}$, hence lightlike for the conformal Lorentzian structure of $\partial\A^{2,1},$ as seen in Remark \ref{222}. In fact $\lambda_{y_0}$ is the orbit of $(x_0,y_0)$ by the action of $\PSL\times\{\text{Id}\},$ or by the (now free) action of $\text{PSO}(2,\R)\times\{\text{Id}\}.$ Here $\text{PSO}(2,\R)$ corresponds to a 1-parameter elliptic subgroup in $\PSL$. In short: 
\[
    \lambda_{y_0}=\PSL\cdot(x_0,y_0)=\text{PSO}(2,\R)\cdot(x_0,y_0).
\]

We refer to $\lambda_{y_0}$ as the \textit{left ruling} through $(x_0,y_0)$, and similarly the \textit{right ruling} is 
\[
    \mu_{x_0}\coloneqq\{(x_0,y)|y\in\R\text{P}^1\}
\]

Our interest in $\partial\A^{2,1}$ is mainly justified by the following: when we will deal with earthquake theory we will often consider $\varphi:\S\to\S$ an orientation-preserving homeomorphism of the circle. The associated graph $\Lambda_\varphi$ via means of the identification given by $\delta$ is a subset of $\partial\A^{2,1}$. We observe that from equivariance of $\delta$ the following holds for every $(\alpha,\beta)\in\PSL\times\PSL$: 
\begin{equation}\label{graphequivariancy}
(\alpha,\beta)\cdot\Lambda_\varphi=\Lambda_{\beta\circ\varphi\circ\alpha^{-1}}.
\end{equation}
We remark one last time that we will consider $\partial\A^{2,1}$ as always implicitly identified with $\T$ via $\delta$.

\section{Geodesic in $\PSL$} We have already seen geodesics in the general Anti-de Sitter space, we would like to specialize here using the model of $\PSL$ and tools from general Lie Group theory. In particular we recall the following (the necessary tools in Lie groups theory are introduced in \cite{bonsanteseppi}): 
\begin{lemma}\label{331}
    Given left-invariant vector fields $V$ and $W$ on $G$, the Levi-Civita connection of a bi-invariant metric has the expression: 
    \[
        \nabla_VW=\frac{1}{2}[V,W].
    \]
    In particular, the Lie group exponential coincides with the pseudo-Riemannian exponential map.
\end{lemma}




We would like to start by considering geodesics through the identity. The Lie algebra of $\PSL$ is isometrically identified with Minkowski space, as seen in Remark \ref{311}, where under such an isometry the stabilizer of a point (namely $\PSL$ acting by means of the adjoint action) corresponds to the group of linear isometries of Minkowski space. In short, this means that we shall distinguish geodesics by their type (timelike, spacelike, lightlike) and those will be equivalent under this action. Moreover, by Lemma \ref{331} it suffices to understand the one-parameter group for the Lie group structure of $\PSL$. We immediately get the following:

\begin{itemize}
    \item Timelike geodesics are, up to conjugacy, of the form: 
    \[
    \begin{bmatrix}
    \cos(t) & -\sin(t) \\
    \sin(t) & \cos(t)
     \end{bmatrix}
\]
namely, under the identification of $\PSL$ with $\text{Isom}(\H^2)$, they are elliptic one-parameter groups fixing a point in $\H^2$. In this example, the tangent vector is the matrix
\[
    \begin{bmatrix}
    0 & -1 \\
    1 & 0
    \end{bmatrix}.
    \]
These are closed geodesics, parametrized by arclenght, of total length $\pi$. 
\item Spacelike geodesics are, again up to conjugacy, of the form: 
\[
    \begin{bmatrix}
    \cosh(t) & \sinh(t) \\
    \sinh(t) & \cosh(t) \end{bmatrix}
\]
with initial velocity: 
\[
    \begin{bmatrix}
        0 & 1 \\
        1 & 0
    \end{bmatrix}.
\]
In the hyperbolic settings, these are loxodromic one-parameter groups, fixing two points in the boundary of $\H^2$ (in this particular case, $\pm 1$). 
\item Finally, lightlike geodesics are the parabolic one-parameter groups conjugate to: 
\[
    \begin{bmatrix}
        1 & t \\
        0 & 1 
    \end{bmatrix},
\]
whose initial vector has indeed zero length.
\end{itemize}

\subsection{Timelike geodesics}
When looking for a complete description of timelike geodesics it suffices to let (the identity component of) the isometry group of $\A^{2,1}$ (namely $\PSL\times\PSL)$ act on $\PSL$ via left and right multiplication. In particular, we can describe the whole space of timelike geodesics of $\A^{2,1}$ as follows: 

\begin{proposition}\label{352}
    There is a homeomorphism between the space of (unparametrized) timelike geodesics of $\A^{2,1}$ and $\H^2\times\H^2$. The homomorphism is equivariant for the action of $\text{Isom}_{0} (\A^{2,1})\simeq \PSL\times\PSL$.
\end{proposition}
\begin{proof}
    The homeomorphism is defined as follows. Given a point $(p,q)\in \H^2\times\H^2,$ we associate to it the subset: 
    \[
        L_{p,q}=\{X\in\PSL\;|\;X\cdot q=p\}
    \]
    By the previous discussion, timelike geodesics through the identity are precisely of the form $L_{p,p}$ for some $p\in\H^2.$ The map $(p,q)\mapsto L_{p,q}$ is equivariant for the natural actions of $\PSL\times \PSL$, namely $(A,B)\cdot L_{p,q}=L_{A\cdot p,B\cdot q},$ which also implies that $L_{p,q}$ is an unparameterized timelike geodesic and that all the unparameterized timelike geodesics are of this form, namely the map we defined is surjective. It remains to show injectivity; if $L_{p,q}=L_{p^{\prime} ,q^{\prime} }$ for $(p,q)\neq(p^{\prime},q^{\prime} )$ then in particular there exists an isometry $X_1$ of $\H^2$ sending $p$ to $q$ and $p^{\prime} $ to $q^{\prime}$, but such an isometry is necessarily unique. Suppose the existence of an $X_2\neq X_1$ isometry of $\H^2$ with the same property, then $X_2^{-1}\circ X_1$ fixes $p,p^{\prime}$, an absurd since the identity is the only isometry of $\H^2$ fixing two different points.
\end{proof}

\subsection{Spacelike geodesics}
Let us consider $\ell$ an oriented geodesics of the hyperbolic plane $\H^2$. From what we have already discussed the one-parameter group of loxodromic transformation fixing $\ell$ as an \textit{oriented} geodesic constitutes a spacelike geodesic through the origin. By an argument analogous to the one given in Proposition \ref{352}, relying on the equivariance of the construction by the actions of $\PSL\times\PSL$, one proves that every spacelike geodesic is of the form: 
\[
    L_{\ell,\jmath}=\{X\in\PSL\;|\;X\cdot\jmath=\ell\;\text{as oriented geodesics}\},
\]
where $\ell$ and $\jmath$ denote oriented geodesics of $\H^2$. We want to emphasize that every (unparameterized, unoriented) spacelike geodesic can be expressed in the above form in two ways, as we could change the orientation of both $\ell$ and $\jmath$. Rephrasing we can state: 
\begin{proposition}\label{353}
    There is a homeomorphism between the space of (unparametrized) oriented spacelike geodesics of $\A^{2,1}$ and the product of two copies of $\partial\H^2\times\partial\H^2\setminus\Delta,$ the space of oriented geodesics of $\H^2$. The homeomorphism is equivariant for the action of $\text{Isom}_0(\A^{2,1})\simeq\PSL\times\PSL$.
\end{proposition}

However we will not be really interested in oriented geodesics, hence we will have identification $L_{\ell,\jmath}=L_{\ell^{\prime},\jmath^{\prime}}$ where with $\ell^{\prime} $ we denote $\ell$ endowed with the opposite orientation.\\
Given a spacelike geodesic, there is a natural notion of \textit{dual} spacelike geodesic, which we define using the projectivity duality between points and planes discussed in section \ref{polarsec}: 
\begin{definition}
    Given a spacelike geodesic $L_{\ell,\jmath}$ in $\A^{2,1}$, the \textit{dual spacelike geodesic} is the intersection of all the spacelike planes dual to points of $L_{\ell,\jmath}.$  
\end{definition}

The construction of the dual geodesic is an involution, let us see it with an explicit example. Consider the geodesic $L_{\ell,\ell}$ through the origin, which consists of the one-parameter loxodromic group of $\PSL$ translating along $\ell$, it can be checked that the dual geodesic consists of all elliptic order-two elements (hence contained in $P_{\text{Id}}$, as we could expect from the definition) whose fixed point lies in $\ell$. To see this we can suppose (as a consequence of the transitivity of the action of $\PSL$ on geodesic of $\H^2$ ) that $\ell$ is the imaginary axis. In this case $L_{\ell,\ell}=\{M_k|k\in\R\}$ where \[
    M_k=\begin{bmatrix}
        e^{-k/2} & 0 \\
        0 & e^{k/2} \\
    \end{bmatrix}.
    \]
It follows from \refeq{geoplanes} that $X\in P_{[M_k]}$ if and only if $\langle X, M_k\rangle=0$, which given a generic matrix
\[
    X=\begin{bmatrix}
        a & b \\

        c & d
        \end{bmatrix}
\]

can be restated as the condition \begin{equation}\label{ellittichine} 
    ae^{\frac{k}{2}}+de^{\frac{-k}{2}}=0.
\end{equation}
As we are asking that Equation \refeq{ellittichine} holds for every $k$ it must be $a=d=0$, hence $X$ (seen as a Möbius transformation) is of the form $X(z)=-\frac{b^2}{z}$ an elliptic transformation of order two with fixed point $ib\in \H^2$.\\
In other words, the dual spacelike geodesic of $L_{\ell,\ell}$ is $L_{\ell,\ell^{\prime}}.$\\
We can explicitly describe the points at infinity in $\partial\A^{2,1}$ of these geodesics. Using Lemma \ref{convergenza}, if $x$ and $y$ are the endpoints at infinity of $\ell$ in $\partial\H^2,$ then any sequence diverging towards an end of $L_{\ell,\ell^{\prime}}\subset \PSL$ maps an interior point towards $x$, and the sequence of inverses towards $y$ (up to switching the two points). In other words, under the identification given by $\delta$ between $\partial\A^{2,1}\simeq\T$, the endpoints of $L_{\ell,\ell}$ are $(x,y)$ and $(y,x)$. A similar argument applied to the geodesic $L_{\ell,\ell^{\prime} },$ which consists of order-two elliptic isometries with fixed point in $\ell$ shows that its endpoints are $(x,x)$ and $(y,y).$\\
Recall the description of the left and right rulings of $\partial\A^{2,1}$ given in \refeq{ruling}, we can conclude that the endpoints of a spacelike geodesic and its dual are mutually connected by lightlike segments in $\partial\A^{2,1}$. \textcolor{red}{Qua andrebbero inserite figure per referenza.}\\
The naturality of the above construction with respect of the action of $\PSL\times\PSL$ let us state: 
\begin{proposition}
    Given a spacelike geodesic $L_{\ell,\jmath}$ of $\A^{2,1}$, its endpoints in $\partial \A^{2.1}$ are $(x_1,y_2)$ and $(y_1,x_2)$, where $x_1$ and $y_1$ are the final and initial endpoints of $\ell$ in $\partial\H^2$, and $x_2$ and $y_2$ are the final and initial endpoints of $\jmath$. The dual geodesic is $L_{\ell,\jmath^{\prime}}$ and has endpoints $(x_1,x_2)$ and $y_1,y_2.$
\end{proposition}


\section{Spacelike planes}
Now we want to study totally geodesic spacelike planes in $\A^{2,1}.$ They are all obtained as the intersection of $\A^{2,1}$ with a projective subspace in the projective space $\text{P}\mathcal{M}(2,\R).$ Hence they are all of the the following form: 
\begin{equation}\label{geoplanes}
    P_{[A]}=\{[X]\in\PSL\;|\;\langle X,A\rangle=0\}
\end{equation}
for some non-zero matrix $A$. The notation is justified by the observation that the plane defined by $P_A$ depends only on the projective class of $A$. The totally geodesic plane is spacelike if and only if $q(A)=-\det A$ is negative. We will call such a plane \textit{dual plane} of $A$, in particular the dual plane $P_\gamma$  of an element $\gamma\in\PSL$ is a spacelike totally geodesic plane. \\
\begin{example}\label{spacelikeisometry}
Before the general treatment we want to focus on a concrete example. Consider $\gamma=\text{Id}\in\PSL.$ Now because of Equation \refeq{quadratic}, $P_{\text{Id}}$ is the subset of $\PSL$ consisting of projective class of unit matrices $X$ with $\text{tr}(X)=0$. By the Cayley-Hamilton theorem, $X^2=-\text{Id},$ hence the elements of $P_{\text{Id}}$ are order-two isometries of $\H^2,$ that is, elliptic elements with rotation angle $\pi$.  We can also observe that $P_{\text{Id}}$ is invariant under the actions of $\PSL$ by conjugation, which corresponds to the diagonal in the isometry group $\PSL\times\PSL$ of $\A^{2,1}$. Using Lemma \ref{convergenza}, we can see that the boundary of $P_{\text{Id}}$ in $\partial\A^{2,1}\simeq \R\text{P}^1\times\R\text{P}^1$ is the diagonal; more precisely: 
\begin{equation}\label{graphino}
    \partial P_{\text{Id}}=\text{graph(Id)}\subset \T.
\end{equation}

Now consider a point $z\in\H^2$, and let us denote by $\mathcal{R}_z$ the order-two elliptic isometry with fixed point $z$. We claim that the map 
\[
    \iota:\H^2\to P_{\text{Id}},\;\;\;\ \iota(z)=\mathcal{R}_z
\]

is an isometry with respect to the hyperbolic metric of $\H^2$ and the induced metric on $P_{\text{Id}}.$ First, the inverse of $\iota$ is simply the fixed-point map $\text{Fix}:P_{\text{Id}}\to\H^2$ sending an elliptic isometry to its fixed point, which also shows that $\iota$ is equivariant with respect to the action of $\PSL$ on $\H^2$ by homographies and on $P_\text{Id}$ by conjugation, since $\text{Fix}(\alpha\gamma\alpha^{-1})=\alpha\text{Fix}(\gamma).$ That is, the following holds: 
\begin{equation}
    \iota(\alpha\cdot p)=\alpha\circ\iota(p)\circ\alpha^{-1}.
\end{equation} 
A direct consequence is that $\iota$ is isometric, since the pull-back of the metric of $P_\text{Id}$ is necessarily $\PSL$-invariant and has constant curvature -1, hence it coincides with the standard hyperbolic metric on the upper half-plane.
\end{example}
This simple example is actually the key case to understand general spacelike totally geodesic planes as every spacelike totally geodesic plane is of the form $P_\gamma$ for some $\gamma\in\PSL.$ To see this, observe that the action of the isometry group of $\A^{2,1}$ on spacelike totally geodetic planes is transitive, and that $P_\gamma=(\gamma,\text{Id})\cdot P_{\text{Id}}$ as the isometry $(\gamma,\text{Id})$ maps $\text{Id}\to\gamma$, and therefore the dual plane of $\text{Id}$ to the dual plane of $\gamma.$ In view of the observations given in \refeq{graphequivariancy} and \refeq{graphino}, we can conclude the following: 
\begin{lemma}\label{32}
    Every spacelike totally geodesic plane of $\A^{2,1}$ is of the form $P_\gamma$ for some orientation-preserving isometry $\gamma\in\PSL,$ and 
    \[
        \partial P_\gamma=\text{graph}(\gamma^{-1})\subset\T.
    \]
\end{lemma}
\section{Timelike planes}
Let us now consider a matrix $A\in\mathcal{M}(2,\R)$ such that $\det(A)=-1.$ The corresponding plane $P_A$ defined by equation \refeq{geoplanes} is a timelike totally geodesic plane. Associated with $[A]$ is an orientation-reversing isometry $\eta$ of $\H^2.$ Moreover, the action of $A$ by homography on $\mathbb{C}\text{P}^1$ preserves $\S$ and switches the two connected components of the complement. The matrix $A$ thus induces an orientation-reversing isometry, up to identifying the two components via $z\to\overline{z}$. We will thus denote $P_{[A]}$ by $P_\eta$. \\
The totally geodesic timelike plane $P_\eta$ can now be parametrized as follows. We have a map: 
\begin{equation}\label{refspa}
    \mathcal{I}\to\mathcal{I}\eta
\end{equation}
from the spaces of reflections $\mathcal{I}$ along geodesic of $\H^2$, with values in $\PSL\simeq\A^{2,1}$. Its image it precisely $P_\eta.$ As seen in the proof of \ref{invertime} we have that an $X$ with determinant $-1$ is an inversion if and only if traceless. As $\det(A)=-1$ we have $\text{adj}(A)=-A^{-1}$, therefore $\langle XA,A\rangle=0$ if and only if $\text{tr}(X)=0$, that is $X$ is traceless hence an involution. This shows that the image of the map defined in Equation \refeq{refspa} is the entire plane $P_\eta$.\\  
In similar fashion to the spacelike case, using the transitivity of the group of isometries on timelike planes, every timelike planes is of the form above. We can show the following: 
\begin{lemma}\label{33}
    Every timelike totally geodesic plane of $\A^{2,1}$ is of the form $P_\eta$ for some orientation-reversing isometry $\eta\in\PSL,$ and 
    \[
        \partial P_\eta=\text{graph}(\eta^{-1})\subset\T.
    \]
\end{lemma}
\begin{proof}
    We want to use Lemma \ref{convergenza} and the parametrization given in Equation \refeq{refspa}. Suppose we have a sequence $\mathcal{I}_n$ such that $\mathcal{I}_n\eta(z)\to x\in\partial\H^2$, for any $z\in\H^2$. Then, using that $\mathcal{I}_n$ is an involution and the continuity of the action of $\eta$ on $\overline{\H}^2$, $(\mathcal{I}_n\eta)^{-1}(z)=\eta^{-1}\mathcal{I}_n^{-1}(z)=\eta^{-1}\mathcal{I}_n(z)=\eta^{-1}(x).$ 
\end{proof}
\begin{observation}
    Since reflections of $\H^2$ are uniquely determined by (unoriented) geodesics, we can consider the map defined in \refeq{refspa} as a map from the space $\mathcal{G}(\H^2)$ of unoriented geodesics of $\H^2$ to $\PSL.$ This map is isometric with respect to the natural metric on $\mathcal{G}(\H^2)$ which makes it identified with the two-dimensional space $\A^{1,1}$ as treated more in detail by El Emam e Seppi in \cite{emam2022gauss} (Example $6.1$).
    \end{observation}
\section{Lightlike planes}
We are only left with case of lightlike totally geodesic planes. Those are the form $P_{[A]}$ for a nonzero matrix $A$ with $\det(A)=0$. Before giving an explicit description of such planes, we want to observe that their boundary will not be a graph in $\T$, unlike the case of spacelike and timelike planes. 

\begin{lemma}\label{35}
    Every lightlike totally geodesic plane of $\A^{2,1}$ is of the form $P_{[A]}$ for some rank one matrix $A$, and: 
    \[
        \partial P_{[A]}=(\text{Im}(A)\times\S)\cup(\S\times\text{Ker}(A)).
    \]
\end{lemma}
\begin{proof}
    The point in $\partial P_{[A]}$ are projective class of rank one matrices satisfying $\langle X,A\rangle=0$, that is, such that $\text{tr}(X\text{adj}(A))=0.$ Given that $X\text{adj}(A)$ has vanishing determinant, by the Cayley-Hamilton theorem $X\text{adj}(A)$ is traceless if and only if it is nilpotent, that is, if and only if $X\text{adj}(A)X\text{adj}(A)=0.$ Now, given that image and kernel of both $X$ and $\text{adj}(A)$ have dimension 1, this happens if and only if: 
    \begin{equation}
        \text{Im}(X)=\text{Ker}(\text{adj}(A))\;\text{or}\;\text{Im}(\text{adj}(A))=\text{Ker}(A).
    \end{equation}
    Now, since $\det(A)=0$ implies $\text{adj}(A)A=A\text{adj}(A)=0$, the relations $\text{Ker}(\text{adj}(A))=\text{Im}(A)$ and $\text{Im}(\text{adj}(A))=\text{Ker}(A)$ hold. Hence $X\in P_{[A]}$ if and only if $\text{Im}(X)=\text{Im}(A)$ or $\text{Ker}(X)=\text{Ker}(A),$ which concludes because of the definition of the homeomorphism $\delta$.  
\end{proof}

Geometrically $\partial P_{[A]}$ is the union of two circles in $\T,$ one horizontal and one vertical, which intersect exactly at the point in $\T$ corresponding to $[A]\in \partial\A^{2,1}$ via $\delta$.\\


%\include{conclusione}
%\appendix
%\include{appendice}
\printbibliography
\end{document}
